\section{Workloads and Setup}
\label{sec:sources}
We use \textbf{Mooncake}~\cite{qin2024mooncake}, replayed production traces from a general LLM serving
stack (conversation, 12{,}031 requests; toolagent, 23{,}608; 512-token blocks), and its third released
workload, synthetic (public datasets, Poisson arrivals), as a control rather than pooled. Alibaba's four
Qwen serving traces~\cite{wang2025kvcachewild} (16-token blocks, full length) are an external
replication, and \textbf{kv-cache-tester}, anonymized agentic coding conversations (64-token blocks), a
second source for reactive caching. Each request lists the hash ids of the blocks it reads, so a
repeated hash is a genuine reuse. Block ids in kv-cache-tester are conversation-scoped, so we namespace
them before pooling; not doing so invents cross-conversation reuse. \Cref{tab:regimes} summarizes the two
primary sources. The third-party traces are not redistributed; the artifact downloads them and checks
pinned SHA-256 digests.

Every static and replayed result is a deterministic CPU computation. Block moves were timed on one host
with an NVIDIA A100 80GB on PCIe Gen4 x16, an AMD EPYC 7513 host, and a rotational RAID-backed disk, on an
otherwise idle GPU while the host's other GPUs ran other jobs; six runs of a 64\,MiB block (an 8B-class
model with grouped-query attention, 512 tokens). The host has no NVMe and no RDMA fabric.

\begin{table}[t]
  \centering \small \setlength{\tabcolsep}{4pt}
  \begin{tabular}{@{}lrr@{}}
    \toprule
     & \makecell{Mooncake\\(serving)} & \makecell{kv-cache-tester\\(agentic)} \\
    \midrule
    requests analysed              & 12{,}031 / 23{,}608 & 4{,}215 (40 conv.) \\
    blocks touched exactly once    & 75.8 / 78.5\% & 15\% \\
    mean reuse per block           & 1.6 / 2.2 & 26 \\
    blocks per request             & 24 / 17 & $\sim$1{,}900 \\
    block size (tokens)            & 512    & 64 \\
    \midrule
    headroom, heuristic vs.\ oracle & 93\%  & 68--73\% \\
    \bottomrule
  \end{tabular}
  \caption{The two sources, at the scales actually analysed. Mooncake pairs are conversation /
  toolagent; its headroom averages the two (93.5\% and 91.8\%). The agentic column is the first 40
  conversation files, except headroom, which is the mean over four randomized draws of 20 and of 40
  conversations (\Cref{sec:sampling}). Its descriptive statistics are sample-dependent: the one-time
  fraction ranges 7.5--21\% over 20--200 conversations. An earlier revision took this column's headroom
  from nested-prefix samples (81--94\%), the protocol \Cref{sec:sampling} retracts.}
  \label{tab:regimes}
\end{table}

